\documentclass[11pt,letterpaper]{article}

\usepackage{graphicx}           % Graphics package
\usepackage{amsmath} 
\usepackage{mathrsfs}
\usepackage[colorlinks=true]{hyperref} 
\usepackage{color}
\usepackage{amsfonts}
\usepackage[textheight=9in, textwidth=7.5in, letterpaper]{geometry}
\usepackage[makeroom]{cancel}
\usepackage{cite}
\usepackage{sectsty}
\usepackage{empheq}
\usepackage{appendix}
\usepackage{enumerate} 
\usepackage{simpler-wick}
\usepackage{amssymb}

\sectionfont{\fontsize{12}{12}\selectfont}
\subsectionfont{\fontsize{12}{12}\selectfont}

\newcommand*\widefbox[1]{\fbox{\hspace{2em}#1\hspace{2em}}}
\newcommand{\LP}{\left(}
\newcommand{\RP}{\right)}
\newcommand{\mc}[1]{\mathcal{#1}}
\newcommand{\R}{\mathbb{R}}
\newcommand{\C}{\mathbb{C}}
\newcommand{\Z}{\mathbb{Z}}
\newcommand{\D}{\partial}
\newcommand{\KET}[1]{\left| #1 \right\rangle }
\newcommand{\BRA}[1]{\left\langle #1 \right| }
\newcommand{\IP}[2]{\left\langle #1 \middle| #2 \right\rangle}
\newcommand{\PA}[2]{\frac{\partial #1}{\partial #2}}



%%%%%%%%%%%%%%%%---------------------MOST USEFUL COMMAND EVER---------------------%%%%%%%%%%%%%%%%%%%%%%
\newcommand{\fixme}[1]{{\bf {\color{red}[#1]}}}
\newcommand{\draftmode}{\usepackage[notref,notcite]{showkeys}}
\providecommand*\showkeyslabelformat[1]{\normalfont\sffamily\footnotesize#1}

%%%%%%%%%%%%%%%


\setcounter{tocdepth}{4}
\setcounter{secnumdepth}{4}

\begin{document}

\title{Quantum Information}

\author{Mathew Calkins}

\date{\today}

 
\maketitle

\abstract{Extremely rough notes on quantum information. Maybe these will turn into something coherent, maybe not!}

\tableofcontents




\section{POVM Measurements}

\subsection{Pure States and Projective Measurements}
As typically introduced, the quantum mechanics of pure states $\KET{\psi}$ living in a Hilbert space $\mc{H}$ connects to experiments via projective measurements. Formally, a projective measurement is described by an indexed collection of projectors \begin{equation}
\{ \Pi_x: \mc{H} \to \mc{H} \ | \ x \in S \}, \quad S \subset \R
\end{equation}
which add up\footnote{Here ``addition" is meant generously to include integration.} to the identity \begin{equation}
\sum_{x \in S} \Pi_x = \mathbf{1}
\end{equation}
so that they equivalently specify an internal direct sum decomposition \begin{equation}
\mc{H} = \bigoplus_{x \in S} \mc{H}_s, \quad \mc{H}_s \equiv \Pi_x (\mc{H}).
\end{equation}
A projective measurement then randomly projects any prepared pure state $\KET{\psi} \in \mc{H}$ \begin{equation}
\KET{\psi} \xmapsto{\mbox{\tiny{proj}}} \frac{\Pi_x \KET{\psi}}{\lVert \Pi_x \KET{\psi} \rVert} = \frac{\Pi_x \KET{\psi}}{\sqrt{\BRA{\psi} \Pi_x \KET{\psi}}}
\end{equation}
and the corresponding $x \in S$ is called the measurement outcome. Assuming WLOG that $\KET{\psi}$ is properly normalized $\IP{\psi}{\psi} = 1$, the probability of any particular outcome is given by the Born rule \begin{equation}
\mbox{Prob}_x = \lVert \Pi_x \KET{\psi} \rVert ^2 = \BRA{\psi} \Pi_x \KET{\psi}.
\end{equation}
If $S \subset \R$ is anything more analytically interesting than a finite set, this object should be interpreted with care as a probability measure. For the present discussion, it suffices to imagine that $S$ is set.


\subsection{Entangled States and POVMs}
The above is all well and good, but the following simple and natural thought experiment suggests a more general kind of process which does not quite take the form of a projective measurement, but which intuitively warrants being called a ``measurement" of our system. Suppose we are interested in the physics of some Hilbert space $\mc{H}_A$, but that there exists a second space $\mc{H}_B$ and that we have the experimental power to perform general joint measurements. A general joint state \begin{equation}
\KET{\Psi}_{AB} \in \mc{H}_A \otimes \mc{H}_B
\end{equation}
is entangled, but for our purposes it will suffice to consider a pure state \begin{equation}
\KET{\Psi}_{AB} = \KET{\psi} \otimes \KET{\phi_{\mbox{\tiny{aux}}}}
\end{equation}
where we think of the second factor as an auxiliary state. Now suppose we perform a projective measurement on the joint state, described as above by a set of projectors \begin{equation}
\{ \Pi_{AB,x} \ | \ x \in S \subset \R\}.
\end{equation}
The Born rule applied to the joint state gives the probability of outcome $x \in S$ as \begin{equation}
\mbox{Prob}_x = \LP \BRA{\psi} \otimes \BRA{\phi_{\mbox{\tiny{aux}}}} \RP  \Pi_{AB,x} \LP \KET{\psi} \otimes \KET{\phi_{\mbox{\tiny{aux}}}} \RP.
\end{equation}
If we think of this as a function of $\KET{\psi}$ at fixed $\KET{\Phi_{\mbox{\tiny{aux}}}}$ and $\Pi_{AB,x}$, we can naturally rewrite the result as the diagonal matrix element of a quadratic form $Q_x$ on $\mc{H}$ with matrix elements \begin{equation}
Q_x: \mc{X}_A \times \mc{H}_A \to \C, \quad \BRA{\phi} Q_x \KET{\chi} \equiv \LP \BRA{\phi} \otimes \BRA{\phi_{\mbox{\tiny{aux}}}} \RP  \Pi_{AB,x} \LP \KET{\chi} \otimes \KET{\phi_{\mbox{\tiny{aux}}}} \RP.
\end{equation}
The RHS is antilinear in the first argument and linear in the second, so this defines a bona fide quadratic form, and we can identify \begin{equation}
\mbox{Prob}_x = \LP \BRA{\psi} \otimes \BRA{\phi_{\mbox{\tiny{aux}}}} \RP  \Pi_{AB,x} \LP \KET{\psi} \otimes \KET{\phi_{\mbox{\tiny{aux}}}} \RP = \BRA{\psi} Q_x \KET{\psi}.
\end{equation}
Recall that the polarization identity on quadratic forms \begin{equation}
Q(u, v) = \frac{Q(u+v, u+v) - Q(u,u) - Q(v,v)}{2}
\end{equation}
tells us a quadratic form's entire action is fixed by its diagonal values, so knowing the value of $\BRA{\psi} Q_x \KET{\psi}$ for all $\KET{\psi} \in \mc{H}_A$ (or more weakly by linearity, just its values on unit-normalized such states) fixes $Q_x$ uniquely.\\
\\
\
The defining properties of $\Pi_{AX,b}$ as a projective measurement, plus judicious use of the polarization identity, imply that each $Q_x$ is a Hermitian positive semi-definite operator\footnote{Here we are slightly redundant, as positive semi-definiteness is only defined for Hermitian operators to begin with.} \begin{equation}
Q_x = Q_x^\dagger, \quad Q_x \geq 0
\end{equation}
and that the collection taken together is a decomposition of the identity \begin{equation}
\sum_{x \in S} Q_x = \mathbf{1}_A.
\end{equation}
These two properties define a \textit{positive operator-valued measurement} (POVM). The reader can convince themselves that any such collection $\{Q_x\}_{x \in S}$ with these two properties can be cooked up as we did here, by a suitable choice of auxiliary state $\KET{\phi_{\mbox{\tiny{aux}}}} \in \mc{H}_B$ and joint projective measurement $\{\Pi_{AB,x} \}_{x \in S}$.



\bibliography{bibfile}
\bibliographystyle{utphys}

\end{document}


